\section{Agent 评测日志归因：从轨迹错误到系统改进}

\subsection{为什么必须看 trajectory log}
在 agent benchmark 中，最终失败通常是多步错误叠加的结果。课程中反复强调，不看轨迹日志就很难判断 ``是模型 reasoning 错误'' 还是 ``工具接口失败''。这也是为什么 reliability 指标必须和 success rate 一起报告。

\begin{knowledgebox}{最小日志字段建议}
\begin{itemize}
    \item 每一步 action 的输入、输出、时间戳；
    \item 工具返回状态码与异常分类；
    \item 中间状态摘要（memory snapshot）；
    \item 终止原因（成功、超时、无效动作、评审拒绝）。
\end{itemize}
\end{knowledgebox}

\subsection{错误归因表：把失败模式量化}
把失败模式结构化后，团队才知道应该优化模型、协议还是 infra。下面的表格示例体现了这一点：表面上 success rate 只有小幅提升，但 invalid action 与 timeout 大幅下降，说明系统稳定性改进显著。

\begin{table}[H]
\centering
\begin{tabular}{p{3.0cm}p{2.5cm}p{2.5cm}p{2.5cm}}
\toprule
\textbf{错误类型} & \textbf{Baseline} & \textbf{New} & \textbf{变化} \\
\midrule
Invalid Action & 12.4\% & 8.1\% & -4.3pp \\
Tool Timeout & 9.7\% & 6.0\% & -3.7pp \\
Wrong Plan & 14.2\% & 13.0\% & -1.2pp \\
Judge Rejection & 6.8\% & 6.1\% & -0.7pp \\
\bottomrule
\end{tabular}
\caption{轨迹级错误归因比单一 success rate 更有诊断价值}
\end{table}

\subsection{一个简化的日志分析脚本}
为了让评测闭环更自动化，可以用轻量脚本聚合 trajectory 失败模式。关键不是脚本复杂度，而是全团队共享同一 taxonomy。

\begin{lstlisting}[caption={聚合 Agent 失败模式的最小示例}]
from collections import Counter

counter = Counter()
for traj in trajectories:
    reason = traj["termination_reason"]
    counter[reason] += 1

total = sum(counter.values())
for k, v in counter.items():
    print(k, f"{v/total:.2%}")
\end{lstlisting}

\begin{warningbox}{没有统一 taxonomy，日志会退化成噪声}
如果每个子团队定义一套 termination reason，跨实验比较会迅速失效，最后又回到 ``只看总分'' 的旧问题。
\end{warningbox}

\subsection{本章小结}
Agent 评测真正的增益来自 ``过程可解释''。轨迹日志、错误归因和统一 taxonomy 是把评测变成可优化系统的必要条件。

